\documentclass{article}

\usepackage[utf8]{inputenc}
\usepackage[T1]{fontenc}
\usepackage{helvet}
\usepackage{geometry}
\usepackage{xcolor}
\usepackage{eso-pic}
\usepackage{fancyhdr}
\usepackage{parskip}
\usepackage{microtype}
\usepackage[english, ukrainian]{babel}
\usepackage{url}
\usepackage{amsmath}
\usepackage{amssymb}
\usepackage{longtable}
\usepackage{booktabs}
\usepackage{hyperref}
\usepackage{titlesec}

\definecolor{nato}{HTML}{293757}
\definecolor{footergray}{gray}{0.65}

\geometry{a4paper,left=3cm,right=3cm,top=3cm,bottom=3cm}
\renewcommand{\familydefault}{\sfdefault}
\setcounter{secnumdepth}{3}

\titleformat{\section}{\color{nato}\Huge\bfseries}{\thesection.}{1em}{}
\titlespacing*{\section}{0pt}{30pt}{12pt}
\titleformat{\subsection}{\color{nato}\LARGE\bfseries}{\thesubsection.}{1em}{}
\titlespacing*{\subsection}{0pt}{20pt}{8pt}
\titleformat{\subsubsection}{\color{nato}\Large\bfseries}{\thesubsubsection.}{1em}{}
\titlespacing*{\subsubsection}{0pt}{10pt}{6pt}

\fancyhf{}
\fancyhead[R]{\small\color{footergray} 29 червня 2026}
\fancyfoot[L]{\small\color{footergray} Технічні вимоги ERP/1. Модуль IAS}
\fancyfoot[R]{\small\color{footergray}\thepage}
\newcommand\HUGE[1]{\fontsize{40}{40}\selectfont #1}
\renewcommand{\headrulewidth}{0pt}
\renewcommand{\footrulewidth}{0.4pt}
\renewcommand{\footrule}{\color{footergray}\hrule width \textwidth height 0.4pt}

\begin{document}
\pagestyle{fancy}

\begin{titlepage}
  \AddToShipoutPictureBG*{\AtPageLowerLeft{\color{nato}\rule{\paperwidth}{\paperheight}}}
  \color{white}\thispagestyle{empty}\vspace*{4cm}\centering
  {\Huge \textbf{ERP/1: Авторизація} \par} \vspace{2cm}
  {\HUGE \textbf{Технічні вимоги} \par} \vspace{2.5cm}
  {\LARGE \textbf{до системи керування цифровими ідентичностями, сертифікатами та політиками безпеки} \par} \vspace{2.5cm}
  \vfill
  {\large 2026 \copyright\ Системи Електронного Урядування \par}
\end{titlepage}

\newpage
\begin{center}{\Huge\color{nato}\bfseries Зміст\par}\end{center}
\vspace{40pt}
\tableofcontents

\begin{abstract}
ERP/1: Служба ідентифікації та авторизації (Identity and Authorization Service --- IAS) є центральним керуючим компонентом безпеки Zencrypted-середовища на базі Erlang/OTP. Поточна реалізація веде стійкий реєстр доменних об'єктів у KVS/Mnesia, підтримує граф зв'язків User--Device--Certificate--Security Profile--VPN Service, валідацію CSR та X.509, покроковий провіженінг пристрою, доставку бажаного стану до VPN через Erlang RPC, ревізії, аудит доставки, реконсиляцію та відновлення після збоїв. Інтеграція з CA у поточній версії містить перехідний CMP/OpenSSL шлях; цільовою є єдина HTTP-взаємодія з CA без залежності IAS runtime від локального OpenSSL. Вимоги щодо BPE/BPMN, EST/CMC, CRL/OCSP, КСЗІ, КЕП та національних криптографічних алгоритмів у цьому документі позначаються окремо як цільові або комплаєнс-вимоги і не трактуються як уже реалізовані функції.
\end{abstract}

\newpage
\section{Умовні скорочення та визначення}
\begin{longtable}{p{4cm}p{10cm}}
\toprule
\textbf{Термін / Скорочення} & \textbf{Значення} \\
\midrule
IAS & Identity and Authorization Service \\
CA & Certificate Authority \\
VPN & Virtual Private Network \\
PKI & Public Key Infrastructure \\
CSR & Certificate Signing Request \\
CMP & Certificate Management Protocol (RFC 4210) \\
EST & Enrollment over Secure Transport (RFC 7030) \\
CMC & Certificate Management over CMS (RFC 5272) \\
BPE & Business Process Engine \\
BPMN & Business Process Model and Notation \\
KVS & Key-Value Storage \\
ABAC & Attribute-Based Access Control \\
FSM & Finite State Machine \\
OTP & Open Telecom Platform \\
\bottomrule
\end{longtable}

\section{Статуси вимог і тверджень}
Щоб не змішувати фактичну реалізацію, план розвитку та нормативні наміри, у документі використовуються такі статуси:
\begin{longtable}{p{3.2cm}p{10.8cm}}
\toprule
\textbf{Статус} & \textbf{Зміст} \\
\midrule
Implemented & Підтверджено поточним кодом і тестами IAS. \\
Partial & Реалізовано частково або лише для конкретного доменного сценарію. \\
Target & Цільова архітектура або функція наступної ітерації. \\
Compliance & Нормативна, сертифікаційна або організаційна вимога; не є доказом реалізації. \\
\bottomrule
\end{longtable}

\section{Загальні відомості}
\subsection{Призначення}
IAS централізує керування цифровими ідентичностями, пристроями, сертифікатами, профілями безпеки, VPN-сервісами та зв'язками між ними. IAS є джерелом бажаного стану для керованих VPN-конфігурацій і координує провіженінг, але не бере участі у пересиланні VPN-трафіку.

\subsection{Поточний функціональний контур}
\begin{itemize}
  \item \textbf{Implemented:} стійке доменне сховище KVS/Mnesia та версіоновані записи.
  \item \textbf{Implemented:} граф User--Device--Certificate--Security Profile--VPN Service.
  \item \textbf{Implemented:} OVPN preview/import, CSR/X.509 validation, certificate lineage та trust evaluation у межах поточного IAS/VPN домену.
  \item \textbf{Implemented:} device-bound provisioning wizard з перевірками переходів, readiness, remediation та відновленням draft.
  \item \textbf{Implemented:} IAS $\rightarrow$ VPN Erlang RPC, канонічні команди, ревізії, digest, idempotence, delivery audit і stale-revision protection.
  \item \textbf{Implemented:} reconciliation, orphan detection, recovery/decommission та restart recovery.
  \item \textbf{Partial:} CA enrollment; поточний CMP/OpenSSL шлях є перехідним і підлягає заміні на HTTP adapter.
  \item \textbf{Target:} виконання оркестрації повного lifecycle через \texttt{synrc/bpmn}/BPE.
\end{itemize}

\subsection{Нормативна рамка}
Застосовні закони, ДСТУ, ISO/IEC/IEEE, NIST та RFC розглядаються як джерела вимог і критеріїв аудиту. Факт згадки стандарту не означає, що сертифікація, атестація або повна технічна відповідність уже завершені.

\section{Архітектурні межі}
\subsection{IAS і CA}
\begin{itemize}
  \item \textbf{AS IS / Partial:} IAS містить CMP enrollment adapter, який запускає локальний OpenSSL CLI і звертається до CA CMP endpoint.
  \item \textbf{Target:} IAS передає CSR та контекст випуску до CA через єдиний HTTP-контракт і отримує випущений сертифікат або структуровану помилку.
  \item Приватний ключ пристрою не повинен передаватися до IAS або CA.
  \item EST і CMC не входять до поточної реалізації IAS і можуть розглядатися лише як майбутні адаптери.
\end{itemize}

\subsection{IAS і VPN}
\begin{itemize}
  \item IAS є джерелом бажаного стану керованих VPN peers і сервісів.
  \item Доставка виконується через distributed Erlang RPC канонічними командами з ревізією та digest.
  \item VPN застосовує команди і зберігає власний durable runtime state.
  \item IAS веде delivery audit, pending/applied стан, виявляє stale revisions та відновлює доставку після restart або reconnect.
  \item Reconciliation порівнює authoritative IAS state з runtime snapshot VPN і формує інциденти та дії recovery/decommission.
\end{itemize}

\subsection{Веб-інтерфейс і сховище}
\begin{itemize}
  \item Browser $\leftrightarrow$ IAS: N2O/NITRO/Cowboy.
  \item IAS persistence: KVS/Mnesia через доменні store-модулі.
  \item \texttt{ias\_demo\_store} використовується як фасад/проєкція для UI та демо-сценаріїв і не визначає окрему продуктивну БД.
\end{itemize}

\section{Функціональні вимоги}
\subsection{Доменна модель}
IAS повинен підтримувати щонайменше такі класи об'єктів:
\begin{itemize}
  \item User, Device, Certificate, SecurityProfile, SecurityPolicy, VPNService;
  \item CertificateEnrollment і CertificateVerification;
  \item ProvisioningTransaction, delivery state та delivery audit;
  \item Relationship graph між користувачами, пристроями, сертифікатами, профілями та сервісами;
  \item Reconciliation incident, orphan recovery/resolution та decommission barrier.
\end{itemize}

\subsection{Provisioning Wizard lifecycle}
Центральний керований процес поточної версії має таку послідовність:
\begin{enumerate}
  \item Scheme.
  \item User.
  \item Device.
  \item Security Profile \& Policy.
  \item VPN Service.
  \item CA Certificate.
  \item Client Certificate.
  \item Relationships.
  \item Material Readiness.
  \item Provisioning.
\end{enumerate}

Поточна реалізація повинна забезпечувати:
\begin{itemize}
  \item стійкий draft із поточним кроком та вибраними доменними ідентифікаторами;
  \item guards переходів і заборону руху вперед без необхідних матеріалів;
  \item автоматичне застосування або ручне виправлення relationship graph;
  \item перевірку certificate material, trust status та authorization decision;
  \item створення idempotent provisioning transaction;
  \item доставку desired state до VPN і відображення результату applied, degraded або failed;
  \item retry, restart recovery і reconciliation після недоступності VPN.
\end{itemize}

\subsection{OVPN import}
OVPN import є допоміжним onboarding-процесом. Система повинна аналізувати конфігурацію, визначати inline CA/certificate/key blocks і параметри підключення, формувати preview та створювати погоджені доменні об'єкти без збереження приватного ключа в IAS.

\subsection{CSR та сертифікати}
\begin{itemize}
  \item IAS повинен приймати зовнішній CSR, перевіряти формат, subject та cryptographic consistency доступними засобами OTP.
  \item IAS повинен зв'язувати випущений сертифікат із CSR, пристроєм, користувачем і профілем безпеки.
  \item Перевірка validity, issuer, fingerprint, public-key continuity та relationship consistency є обов'язковою.
  \item CRL/OCSP online validation є \textbf{Target}, доки окрема інтеграція не реалізована і не протестована.
\end{itemize}

\subsection{Авторизація і довіра}
Поточна реалізація використовує профілі, політики, атрибути доменних об'єктів, relationship graph і certificate trust для рішень у IAS/VPN сценарії. Це \textbf{Partial ABAC}: універсальний policy language, зовнішні attribute sources, геолокаційні правила та загальний-purpose ABAC engine є \textbf{Target}.

\subsection{Provisioning і reconciliation}
\begin{itemize}
  \item Кожна керована зміна повинна мати канонічне представлення, revision і digest.
  \item Повторна доставка однакової команди повинна бути idempotent.
  \item Застарілі ревізії не повинні перезаписувати новіший стан.
  \item Pending і applied heads повинні переживати restart.
  \item Orphan, fingerprint mismatch, stale revision та unauthorized runtime state повинні відображатися як reconciliation incidents.
  \item Recovery/decommission дії повинні бути аудованими та захищеними від повторного застосування старого desired state.
\end{itemize}

\section{Нефункціональні вимоги}
\subsection{Архітектура і надійність}
\begin{itemize}
  \item Реалізація на Erlang/OTP із supervision trees і restart-safe state transitions.
  \item Authoritative persistence у KVS/Mnesia.
  \item Явні schema versions та upgrade/migration logic для стійких записів.
  \item Автоматизовані EUnit і Common Test сценарії для доменних store, provisioning, RPC, reconciliation та recovery.
  \item Кількісні SLA, capacity і latency targets повинні підтверджуватися окремими benchmark/acceptance tests; до цього вони мають статус \textbf{Target}.
\end{itemize}

\subsection{Безпека}
\begin{itemize}
  \item Приватні ключі користувачів і пристроїв не зберігаються в IAS.
  \item Секрети, ключовий матеріал і повні credential values не записуються в audit/log output.
  \item Transport security, node authentication та secret distribution повинні бути визначені для кожного deployment profile.
  \item КСЗІ Г3, CAdES-X Long, ДСТУ 4145-2002, ДСТУ ГОСТ 28147:2009 та at-rest encryption мають статус \textbf{Compliance/Target}, доки немає окремих реалізаційних і атестаційних доказів.
\end{itemize}

\subsection{Логування і аудит}
\begin{itemize}
  \item IAS повинен аудувати lifecycle transitions, issuance/import, authorization decisions, provisioning delivery, reconciliation incidents і recovery actions.
  \item Audit records повинні мати ідентифікатор об'єкта, revision, результат, причину та часову мітку.
  \item WORM storage, ELK/Loki export і довгострокове нормативне зберігання є окремими deployment/compliance вимогами, а не автоматичною властивістю поточної реалізації.
\end{itemize}

\subsection{Документація}
\begin{itemize}
  \item Технічні вимоги, техноробочий проєкт і політики повинні використовувати однакові статуси Implemented, Partial, Target і Compliance.
  \item API/OpenAPI, operator manuals і security profiles створюються відповідно до фактично наявних інтерфейсів.
  \item BPMN/FSM-діаграми повинні спочатку відображати AS IS lifecycle; BPE execution додається після реалізації відповідної runtime integration.
\end{itemize}

\section{Простежуваність вимог}
\begin{longtable}{p{2.3cm}p{5.4cm}p{2.2cm}p{4.1cm}}
\toprule
\textbf{ID} & \textbf{Вимога} & \textbf{Статус} & \textbf{Доказ / зона реалізації} \\
\midrule
IAS-FR-001 & Стійкий реєстр User/Device/Certificate/Profile/VPN Service & Implemented & domain/store modules, KVS/Mnesia tests \\
IAS-FR-002 & CSR і X.509 validation & Implemented & certificate/provisioning modules and tests \\
IAS-FR-003 & CA enrollment через HTTP без OpenSSL у IAS runtime & Target & заміна legacy CMP/OpenSSL adapter \\
IAS-FR-004 & Device-bound provisioning wizard & Implemented & wizard, readiness, relationships, authorization \\
IAS-FR-005 & IAS $\rightarrow$ VPN RPC delivery & Implemented & VPN client/delivery modules and CT suites \\
IAS-FR-006 & Revision, digest, idempotence і stale protection & Implemented & delivery/provisioning state and tests \\
IAS-FR-007 & Reconciliation, orphan recovery і decommission & Implemented & reconciliation/recovery modules and CT suites \\
IAS-FR-008 & BPE/BPMN runtime orchestration & Target & planned migration of wizard orchestration \\
IAS-FR-009 & CRL/OCSP, EST і CMC adapters & Target & absent from current runtime \\
IAS-NFR-001 & КСЗІ Г3, КЕП, ДСТУ crypto, at-rest encryption & Compliance & requires separate implementation and evidence \\
\bottomrule
\end{longtable}

\section{Додаток А. Межі поточної версії}
Поточна версія не повинна описуватися як така, що вже забезпечує повну універсальну ABAC-платформу, EST/CMC, online CRL/OCSP, загальний REST/OpenAPI surface, TLS 1.3 для всіх каналів, RocksDB/NVMe persistence, WORM audit, гарантовані 100000 пристроїв, 50 мс latency, КСЗІ Г3, CAdES-X Long або реалізацію національної криптографії. Такі твердження можуть бути включені лише як Target або Compliance з окремими acceptance criteria.

\section{Додаток Б. Цільова послідовність розвитку}
\begin{enumerate}
  \item Синхронізувати \texttt{ias.tex} і \texttt{ias\_pro.tex} з поточним кодом.
  \item Замінити server-side OpenSSL/CMP у IAS на канонічний HTTP adapter до CA.
  \item Прогнати і стабілізувати повний provisioning wizard lifecycle.
  \item Зафіксувати AS IS FSM/BPMN для реального lifecycle.
  \item Поетапно перенести orchestration на \texttt{synrc/bpmn}/BPE, зберігаючи доменні service modules і тестовану поведінку.
  \item Повторно синхронізувати технічні вимоги, техноробочий проєкт і security profiles.
\end{enumerate}

\end{document}
